% Part: applied-modal-logics
% Chapter: epistemic-logic
% Section: relational-models

\documentclass[../../../include/open-logic-section]{subfiles}

\begin{document}

\olfileid{aml}{el}{rel}

\olsection{النماذج العلائقية}

أصل الدلالة في منطق المعرفة هو أصلها في المنطق الجهوي المعتاد: نموذج
علائقي فيه عوالم، وإسناد يعين عند أي عالم تصدق ال!!{propositional variable}s.
فإن كان الفاعل واحدًا كانت علاقة النفاذ واحدة؛ وإن تعدد الفاعلون تعددت
علاقات النفاذ، فلكل~$\OLId{latin-ordinary}{a}$ من مجموعة الرموز~$\OLId{latin-looped}{G}$ علاقة.

فــ\emph{النموذج العلائقي}، إذن، مجموعة عوالم تصل بينها علاقات نفاذ
ثنائية بعدد الفاعلين، ومعها إسناد يعين ال!!{propositional variable}s
الصادقة عند كل عالم.

\begin{defn}
  \emph{نموذج} لغة المعرفة المتعددة الفاعلين هو الثلاثي
  $\mModel{M}=\tuple{\OLId{latin-looped}{W},(\OLId{latin-looped}{R}_\OLId{latin-ordinary}{a})_{\OLId{latin-ordinary}{a} \in \OLId{latin-looped}{G}},\OLId{latin-looped}{V}}$، وفيه:
  \begin{enumerate}
  \item $\OLId{latin-looped}{W}$ مجموعة غير فارغة من «العوالم»؛
  \item لكل~$\OLId{latin-ordinary}{a} \in \OLId{latin-looped}{G}$، تكون~${\OLId{latin-looped}{R}}_\OLId{latin-ordinary}{a}$ علاقة نفاذ ثنائية على~$\OLId{latin-looped}{W}$؛
  \item $\OLId{latin-looped}{V}$ دالة تسند إلى كل~!!{propositional
    variable}~$\OLId{latin-ordinary}{p}$ مجموعة~$\OLId{latin-looped}{V}(\OLId{latin-ordinary}{p})$ من العوالم الممكنة.
  \end{enumerate}
  فإذا كان~$\OLId{latin-looped}{R}_\OLId{latin-ordinary}{a} ww'$ قلنا إن الفاعل~$\OLId{latin-ordinary}{a}$ \emph{يمكنه النفاذ إلى}~$\OLId{latin-ordinary}{w}'$
  من~$\OLId{latin-ordinary}{w}$؛ وإذا كان~$\OLId{latin-ordinary}{w} \in \OLId{latin-looped}{V}(\OLId{latin-ordinary}{p})$ قلنا إن~$\OLId{latin-ordinary}{p}$ \emph{صادق عند}~$\OLId{latin-ordinary}{w}$.
\end{defn}

فليس وراء ذلك إلا آلية المنطق الجهوي النظامي مع تكثير علاقات النفاذ.
ونحمل علاقة كل فاعل على حاله المعلوماتية: فقولنا~$\OLId{latin-looped}{R}_\OLId{latin-ordinary}{a} ww'$ يعني، في
التفسير المعتاد، أن~$\OLId{latin-ordinary}{w}'$ متسق مع معلومات~$\OLId{latin-ordinary}{a}$ عند~$\OLId{latin-ordinary}{w}$؛ أي إن~$\OLId{latin-ordinary}{a}$ عند
$\OLId{latin-ordinary}{w}$ لا يميز العالم~$\OLId{latin-ordinary}{w}$ من العالم~$\OLId{latin-ordinary}{w}'$.

\end{document}
